\subsection{An explicit higher special cotangent class}

We test the proposed horizontal-support bound on a cellular row-zero
resolution.  The calculation uses the special differential relation
$d\gamma_m(x)=\gamma_{m-1}(x)\,dx$ throughout.  In particular, the
complex below is given explicitly, without identifying its modules
with arbitrary Beck modules.

\paragraph{A-S09-09 (Proved; unbounded-weight classes in row one).}
Let $S=\Gamma_{\mathbb Z}(T)$, let $I_2$ be its full weight tail
starting at two, and put $R=S/I_2$.  There is a graded cellular free
simplicial PD $S$-algebra resolution $Q\to R$ such that, for every
odd prime $p$, the group
\[
 \pi_2\bigl(\Omega^1_{\mathrm{PD}}(Q/S)\bigr)^{2p}
\]
contains an element of exact order $p$.  Its weight-six group is
exactly $\mathbb Z/3$.  Consequently the proposed eventual support
bound $n\leq 2q-1$ for the homology of row $q$ is false for this
construction: it fails in row $q=1$ in unbounded weights.

\begin{proof}
Write $M=\mathbb Zt$, with $\operatorname{wt}(t)=1$, and write
$G(V)=\Gamma^+_{\mathbb Z}(V)$ for the positive part of the free
PD algebra.  Its unit and multiplication are $\eta$ and $\mu$.
The positive part of $R$ is $M$, with structure map
\[
 \alpha:G(M)\longrightarrow M,\qquad
 \alpha(\gamma_1(t))=t,\quad
 \alpha(\gamma_j(t))=0\quad(j>1).
\]
Use the monadic resolution
$Q_n=\mathbb Z\oplus G^{n+1}(M)$, with its standard multiplication,
faces given by $\mu$ and $\alpha$, and degeneracies given by
$\eta$.  The maps from $S=\mathbb Z\oplus G(M)$ send $T$ to
the completely unary iterated unit.  They commute with every
face and degeneracy, so $Q$ is a simplicial PD $S$-algebra.
The unit gives an extra degeneracy after forgetting the algebra
structure.  It is additive and proves that $Q\to R$ is a
weightwise weak equivalence.

Every $G^n(M)$ is free as an abelian group on its divided-power
monomials.  A degeneracy inserts a layer consisting entirely of
unary operations and sends basis elements injectively to basis
elements.  Removing all such layers gives the unique nondegenerate
ancestor of a basis element.  Hence, over $S$, the simplicial
algebra is built by adjoining the remaining free simplex
generators along their prescribed boundaries.  This proves
the asserted cellular property, including $Q_0=S$.

Let
\[
 L=R\otimes_Q\Omega^1_{\mathrm{PD}}(Q/S).
\]
The map $\Omega^1_{\mathrm{PD}}(Q/S)\to L$ is a weak equivalence.
Here is the required comparison.  Filter the differential module
by the cellular generators just described.  A generator of
simplicial dimension $n$ contributes, on the associated quotient,
$Q\otimes_{\mathbb Z}\mathbb Z[\Delta^n/\partial\Delta^n]$
with its indicated weight; the corresponding quotient after base
change has $R$ in place of $Q$.  The quotient comparison is a
weak equivalence because $Q\to R$ is one and the normalized cell
complex is bounded and free over $\mathbb Z$.  The filtration
inclusions split in every simplicial degree.  The long exact
homology sequence, followed by exactness of filtered colimits,
proves the comparison.  This is also an independent check of
the particular comparison used in SA08 C15.

We describe $N(L)$ completely.  Its degree-zero group is zero.
In degree $n>0$, use the quotient of
$R\otimes_{\mathbb Z}G^n(M)$ by the monomials with a wholly
unary layer.  The first face is special differentiation of the
outer monomial, with its coefficient evaluated in $R$; the
interior faces substitute two successive $G$-layers; the last
face applies $\alpha$ to the innermost layer.  The differential
is the alternating sum of these faces.  The last face is zero
on every nondegenerate basis element.  A normalized basis element
of parameter weight $w$ in degree $n$ must have $n$ branching
layers, whence $n\leq w-1$.  Thus every fixed-weight complex in
this description is finite.

Set $g_j=\gamma_j(t)\in G(M)$ and $u=g_2$.  An expression in
the $g_j$ in the following formulas means a monomial in the
\emph{next} free divided-power algebra.  For an odd prime $p$,
put
\[
 c_p=\frac{(2p-3)!!}{p}\in\mathbb Z,
 \qquad
 \xi_p=\gamma_p(u)-c_p\,u\,g_{2p-2}\in N_2(L)^{2p}.
 \tag{S09.11}
\]
The integer $c_p$ is defined because the odd product
$1\cdot3\cdots(2p-3)$ includes the factor $p$.
Both terms are nondegenerate.  The first and last faces vanish:
their positive-weight coefficients evaluate to zero in $R$.
For the middle face, the divided-power composition and product
identities give
\[
 \mu\gamma_p(g_2)=(2p-1)!!\,g_{2p},
 \qquad
 \mu(g_2g_{2p-2})=p(2p-1)\,g_{2p}.
\]
The choice of $c_p$ makes these equal.  Thus $d\xi_p=0$.

In degree three define
\[
 z_p=\eta(u)\,\gamma_{p-1}(u)\in G(G(G(M))).
\]
The two displayed factors are elements of $G(G(M))$, and
their product is taken in the outer free PD algebra.
This is again nondegenerate.  Its first and last faces vanish.
The two interior substitutions give respectively
$p\gamma_p(u)$ and $(2p-3)!!\,u g_{2p-2}$.  With their
alternating signs,
\[
                       dz_p=-p\xi_p.                 \tag{S09.12}
\]

Define $\ell_p:N_2(L)^{2p}\to\mathbb F_p$ by extracting
the coefficient of $\gamma_p(u)$ in the coefficient-one
summand.  We check $\ell_p d=0$ on every normalized degree-three
basis element.  The first face has a positive-weight coefficient
unless its outer layer is unary; the latter elements are
degenerate.  It contributes zero to $\ell_p$.  The last face
is zero as already observed.  For the first interior face, a
term contributing to $\gamma_p(u)$ consists entirely of
intermediate monomials $\gamma_j(u)$.  Apart from the two
degenerate cases, its coefficient is
\[
 \frac{p!}{\prod_j (j!)^{a_j}\prod_j a_j!},
 \qquad \sum_j j a_j=p,
\]
with every $j$ and $a_j$ strictly less than $p$.  It is therefore
divisible by $p$.  For the other interior face, an intermediate
monomial substituting to $u=\gamma_2(t)$ must be one of
$\eta(u)$ and $\gamma_2(\eta(t))$.  The pure $p$-th divided
powers of these are degenerate at one of their two inner layers.
The remaining contributions have coefficients
$\binom pa$, $0<a<p$, which are divisible by $p$.
This exhausts the monomial basis.  Thus $\ell_p d=0$ and
$\ell_p(\xi_p)=1$.  Equations (S09.11)--(S09.12) prove that
$[\xi_p]$ has exact order $p$.

For completeness, the entire weight-six group can be obtained
with a small presentation.  Its normalized ranks in positive
degrees are $2,14,37,41,16$.  In degree one use $(g_6,Tg_5)$.
In degree two take the ordered basis displayed in the first
column of the following table.  Substitution of the thirty-seven
degree-three monomials gives the quotient in the second column,
where $a,b$ are free and $3\tau=0$:
\[
\begin{array}{c|r}
g_1g_2g_3&12a\\
g_1g_5&b\\
\gamma_2(g_1)\gamma_2(g_2)&9a\\
\gamma_2(g_1)g_4&3a\\
\gamma_3(g_1)g_3&4a\\
\gamma_4(g_1)g_2&3a\\
g_2g_4&3a\\
\gamma_3(g_2)&3a+\tau\\
\gamma_2(g_3)&2a\\
Tg_1\gamma_2(g_2)&18a-15b\\
Tg_1g_4&6a-5b\\
T\gamma_2(g_1)g_3&12a-10b\\
T\gamma_3(g_1)g_2&12a-10b\\
Tg_2g_3&12a-10b
\end{array}
\tag{S09.13}
\]
One can check the presentation by eliminating the unit-coefficient
relations first.  Writing the basis as $e_0,\ldots,e_{13}$,
the remaining relations give $3e_8=2e_6$ and
$3(e_7-e_6)=0$; put $a=e_6-e_8$, $b=e_1$,
$\tau=e_7-e_6$.  They express every $e_i$ exactly as in
(S09.13).  Substitution of these expressions into every face
relation gives zero, so the indicated assignments are mutually
inverse presentations, rather than just a surjection onto the
quotient.  The induced differential sends
\[
       a\longmapsto(-5,0),\qquad
       b\longmapsto(-6,1),\qquad \tau\longmapsto0.
\]
Its restriction to the free part is injective.  Its kernel is
therefore exactly $\mathbb Z/3\tau$, which proves the
weight-six assertion.

The independent certificate constructs all the integral normalized
matrices in weights two through six.  It also checks 5,328
unnormalized face identities, the cycle and boundary in
(S09.11)--(S09.12), and the functional $\ell_3$ on all
thirty-seven boundaries.  The degree-three matrix has Smith
factors $1$ eleven times and $3$ once; the degree-two
differential has rank two, agreeing with (S09.13).
\end{proof}

The prime-local range $n<pq$ is not refuted by these classes:
after localization at an odd prime $p$, the class displayed above
has $n=2<p$ and $q=1$.  We have not established a bound in that
range for all cotangent homology, nor a comparison from this
row-zero resolution to an admissible common source.  The statement
proved here is the explicit failure of the proposed integral
support bound, with the exact construction and maps specified.

\paragraph{A-S09-10 (Proved; an exact failure of the prime-local slope).}
For the same resolution $Q\to R=S/I_2$, one has
\[
 \pi_3\bigl(\Omega^1_{\mathrm{PD}}(Q/S)\bigr)^9
       \otimes_{\mathbb Z}\mathbb Z_{(3)}\ne0.
\]
Thus the support inequality $n<pq$ also fails, at
$p=3,n=3,q=1$, for this particular free resolution.  This is a
finite-weight assertion; it does not establish or refute an
eventual prime-local support bound.

\begin{proof}
Use the explicitly described complex $C=N(L)$ from A-S09-09.
Put $a=g_1$, $b=g_2$, and $c=g_3$ in $G(M)$.  Form the following
elements of $G^3(M)$, where the displayed $\eta$ maps insert
the indicated elements of $G^2(M)$ into the outer free algebra:
\[
 \begin{split}
 A&=\gamma_3(ab),\\
 B&=\eta(\gamma_3(a))\,\eta(\gamma_3(b)),\\
 D&=\eta(\eta(c))\,\eta(\gamma_2(c)),\\
 \zeta&=A-6B+9D\in C_3^9.
 \end{split}
\]
All these monomials are nondegenerate.  In $C_2^9$ write
$E=\gamma_3(a)\gamma_3(b)$, $F=\gamma_3(c)$,
and $J=cg_6$.  Their faces give, integrally,
\[
 dA=-6E+27F,\qquad
 dB=-E+15J,\qquad
 dD=-3F+10J.
 \tag{S09.14}
\]
Indeed the two relevant identities for the first expression are
$\gamma_3(ab)=6\gamma_3(a)\gamma_3(b)$ and
$\mu(ab)=3c$.  For the other two, use
$\mu\gamma_3(b)=15g_6$ and
$\mu\gamma_2(c)=10g_6$.  The first and last faces vanish.
Equation (S09.14) proves $d\zeta=0$.

Let $\ell:C_3^9\to\mathbb F_3$ extract the coefficient of $A$.
We verify on the full basis that $\ell d:C_4^9\to\mathbb F_3$
is zero.  The first face has a positive-weight coefficient unless
the outer layer is unary, which is degenerate.  The last face
vanishes.  For the first interior substitution, the prime
multinomial argument of A-S09-09 applies to three copies of $ab$.
Every nondegenerate contribution is divisible by three.

For the middle substitution, a single middle monomial substituting
to $ab$ is either $\eta(ab)$ or $\eta(a)\eta(b)$: the output
contains two distinct variables, each once.  The pure third
power of either preimage has a wholly unary layer.  Mixtures of
the two have coefficient $\binom31$ or $\binom32$.
For the last interior substitution, the preimage of $a$ is uniquely
$\eta(a)$, and the preimages of $b$ are $\eta(b)$ and
$\gamma_2(a)$.  Thus the two possible preimages of $ab$ are
$\eta(a)\eta(b)$ and $\eta(a)\gamma_2(a)$.  A pure third
power has, respectively, a wholly unary middle or innermost
layer.  Mixtures again have a binomial coefficient divisible by
three.  This exhausts the normalized basis and proves $\ell d=0$.
Since $\ell(\zeta)=1$, the integral cycle cannot be a boundary
over $\mathbb Z_{(3)}$: multiplying a proposed boundary by a
denominator prime to three and applying $\ell$ would give a
nonzero element of $\mathbb F_3$ equal to zero.

The finite certificate independently constructs the whole
weight-nine complex through degree five and checks (S09.14)
and all 2,417 degree-four boundary values of $\ell$.  Its
dimensions in degrees one through five are
$(2,48,497,2417,6018)$.  The ranks of $d_2,d_3,d_4,d_5$
over $\mathbb F_3$ are $(2,46,450,1966)$; over
$\mathbb F_{101}$ they are $(2,46,451,1966)$.
These last ranks imply rational vanishing in degrees two
through four by exact inequalities: ranks over $\mathbb Q$
are at least their reductions, while $d^2=0$ bounds each
adjacent rank sum by the dimension of the middle term.
The displayed lower bounds already attain those dimensions.
Thus the rational homology vanishes in these degrees.
Moreover $H_2(C\otimes\mathbb F_3)=0$, whereas
$H_3(C\otimes\mathbb F_3)=\mathbb F_3$.
The universal-coefficient exact sequence then shows that
$H_3(C)$ has nonzero three-primary torsion.  This last
observation supplements the direct cycle proof; it is not
an inference that a single finite-field rank automatically
equals its rational rank.
\end{proof}
